%% fig-chain.tex -- Figure: the governance chain in three registers.
%%
%% Upper: the seven-term abstract chain -- a reference shape, NOT a stage count.
%% Middle: the concrete six-stage prompt refinement path, the fullest instance.
%% Lower: the durable residue each term leaves behind.
%%
%% Register labels are rotated into the left margin so the vertical space between
%% registers stays clear for the correspondence lines, which are deliberately not
%% one-to-one -- that mismatch is the figure's argument.
%%
%% Type scale: \figH headings, \figB detail. Seven columns across 7in leaves each
%% cell about 2.1cm, so detail lines are held to two or three short words.

\begin{figure}[tbp]
\centering
\begin{cbwide}
\begin{tikzpicture}[x=1cm,y=1cm]

  \def\X{1.86}    % content left edge (left of it: rotated register labels)
  \def\W{15.60}   % content width
  \def\G{0.20}

  % ================= upper register: the abstract chain =====================
  \foreach [count=\i from 0] \term/\leaves in {%
      {\mode{Signal}}/{a failure worth\\acting on},
      {\mode{Evidence}}/{traces + test\\sets assembled},
      {\mode{Candidate}}/{a proposed\\better version},
      {\mode{Measure-}\\\mode{ment}}/{scored on that\\evidence},
      {\mode{Decision}}/{an operator\\applies, or not},
      {\mode{Release}}/{the \livetarget{}\\moves},
      {\mode{Trace}}/{the next signal\\arrives}}
  {
    \pgfmathsetmacro{\w}{(\W-6*\G)/7}
    \node[control, anchor=north west, text width=\w cm-8pt, minimum height=1.18cm,
          inner xsep=2pt] (c\i) at ({\X+\i*(\w+\G)}, 0.30)
      {\term\\[0.25ex]{\figB\leaves}};
  }
  \foreach \i [evaluate=\i as \j using int(\i+1)] in {0,...,5}
    { \draw[ctlflow] (c\i.east) -- (c\j.west); }

  % the loop closes, routed above the register so the lane below stays clear
  \draw[feedback] (c6.north) -- ({\X+\W-1.0},0.78) -- (2.96,0.78) -- (c0.north);
  \node[anchor=south, font=\figB, text=cbAsync, inner sep=1.8pt]
    at ({\X+\W/2},0.80)
    {the loop closes --- client code keeps resolving \prodalias{}, so the next
     production trace \emph{is} the next signal};

  % ================= middle register: the concrete path =====================
  \foreach [count=\i from 0] \s/\name/\detail in {%
      muted/{Production trace}/{\mlflow{} trace\\+ assessment},
      human/{1.~Verify}/{\textbf{human} ---\\is it real?},
      control/{2.~Diagnose}/{LLM root cause\\on the evidence},
      control/{3.~Optimize}/{policy-selected\\optimizer},
      control/{4.~Evaluate}/{EvalProvider +\\regression gate},
      human/{5.~Apply}/{\textbf{operator} ---\\diff + scores},
      extern/{6.~Promote}/{audited alias\\rotation}}
  {
    \pgfmathsetmacro{\w}{(\W-6*\G)/7}
    \node[\s, anchor=north west, text width=\w cm-8pt, minimum height=1.18cm,
          inner xsep=2pt] (s\i) at ({\X+\i*(\w+\G)}, -1.96)
      {\textbf{\name}\\[0.25ex]{\figB\detail}};
  }
  \foreach \i [evaluate=\i as \j using int(\i+1)] in {0,...,5}
    { \draw[flow] (s\i.east) -- (s\j.west); }

  % ---- the correspondences, drawn because they are not one-to-one ----------
  \foreach \a/\b in {0/1, 1/2, 2/3, 3/4, 4/5, 5/6, 0/0, 1/1}
    { \draw[weak, densely dotted, draw=cbControl, line width=0.5pt]
        (c\a.south) -- (s\b.north); }

  % ================= lower register: the durable residue ====================
  \foreach [count=\i from 0] \r in {%
      {verification\\item},
      {refinement job\\+ evidence},
      {diagnosis +\\candidate},
      {scores + gate\\decision},
      {apply action +\\provenance anchor},
      {checkpoint +\\audit row},
      {new traces,\\SLO incidents}}
  {
    \pgfmathsetmacro{\w}{(\W-6*\G)/7}
    \node[store, anchor=north west, text width=\w cm-8pt, minimum height=0.86cm,
          font=\figB, inner xsep=2pt] at ({\X+\i*(\w+\G)}, -3.56) {\r};
    \draw[flow, draw=cbStore]
      ({\X+\i*(\w+\G)+\w/2}, -3.18) -- ({\X+\i*(\w+\G)+\w/2}, -3.52);
  }

  % ================= rotated register labels ================================
  % Short labels on purpose: a rotated multi-word label is taller than the register
  % it names, so it runs into the neighbouring one.
  \foreach \y/\lab in {-0.29/{CONCEPTS}, -2.55/{STAGES}, -3.99/{RECORDS}}
  {
    \node[anchor=center, font=\figB\bfseries, text=cbMuted] at (0.74,\y)
      {\rotatebox{90}{\lab}};
  }

\end{tikzpicture}
\end{cbwide}
\caption{The \chainname{} in three registers. The top row is a \emph{reference
shape}, not a worker pipeline: it names seven concepts a refinement-capable
\family may occupy. The middle row is the concrete prompt path --- the fullest
realization of that shape --- with six numbered stages and exactly two human
decisions. The dotted correspondences are deliberately not one-to-one, and the
mismatch is the figure's argument: \cEvidence{} has no stage of its own, because
evidence assembly is folded into the transition from \mode{Verify} toward
\mode{Diagnose}; \cSignal{} spans both an incoming production trace and the human
confirmation that the failure is real; and \cTrace{} is simultaneously the
seventh term and the input to the next iteration. Reading the seven terms as
seven worker stages is the single most common misreading of this architecture.
Other families reuse the shape without inheriting its guarantees: a workflow is
measured by manifest replay, a tool by a revision-fenced deterministic suite, a
knowledge base by retrieval calibration (\tabref{tab:families}). The bottom row
is what separates this from an observability dashboard --- every term deposits
durable state, so a release is reconstructable from records rather than from
recollection.}
\label{fig:chain}
\figkey
\end{figure}
